\subsection{معرفی و انگیزه}
این ماژول در چارچوب پژوهش حاضر پیاده‌سازی و ارزیابی شده و فاقد مقاله مستقل جداگانه در ادبیات تحقیق است.

کوانتایزاسیون وزن‌ها و فعال‌سازی‌ها به قالب‌های با بیت‌پایین (مانند اعداد صحیح ۸ بیتی، \lr{INT8}) یکی از پرکاربردترین استراتژی‌های کاهش حجم مدل و ارتقای پهنای باند حافظه در استنتاج شبکه‌های زبانی است. با این حال، کوانتایزاسیون استاندارد و یکنواخت به شدت در برابر ناهمگونی دامنه‌های ویژگی آسیب‌پذیر است؛ در لایه‌های ترنسفورمر، برخی کانال‌های فعال‌سازی دارای مقادیر دورافتاده (\lr{Outliers}) با دامنه‌های بسیار بزرگ هستند که در صورت نادیده گرفته شدن، خطای گردکردن سایر ویژگی‌های ظریف را تشدید کرده و به تخریب چشمگیر خروجی شبکه می‌انجامند.

ماژول فشرده‌ساز \lr{FAQQuantizer} با نام‌های ثبتی \lr{faq} و \lr{faq\_quantizer} («کوانتایزه‌کننده آگاه از فعال‌سازی ویژگی‌ها» یا \lr{Feature-Activation-Aware Quantization})، بر مبنای مقیاس‌گذاری متقابل کانال‌ها عمل می‌کند. این روش با پایش آماری ممان دوم فعال‌سازی‌ها در فاز کالیبراسیون، گام‌های کوانتایزاسیون را در هر کانال به گونه‌ای مقیاس‌بندی می‌کند که نویز ناشی از خطای کوانتایزاسیون در جهت کانال‌های با بار اطلاعاتی بیشتر کمینه گردد. این ماژول همچنین با پیاده‌سازی لایه اختصاصی \lr{FAQLinear} و تقریب گرادیان سرراست (\lr{Straight-Through Estimator, STE})، بستر مناسبی برای تنظیم دقیق با \lr{LoRA} فراهم می‌سازد.
